\section{Asymptotic rank typicality without isotropy}
\label{sec:asymptotic-rank-typicality}

This section separates a mathematical random-orientation statement from the empirical question of whether a physical circuit/readout pair realizes that orientation regime.  Let $C\succeq0$ be an arbitrary real $N\times N$ covariance operator with $\Tr C=1$, and let $P$ be a rank-$r$ orthogonal projector drawn uniformly from the real Grassmannian, independently of $C$.  Define
\begin{equation}
 \rho(P,C)=\frac{\Tr(PC)}{r/N},
 \qquad
 d_{\rm eff}(C)=\frac{1}{\Tr(C^2)}.
\end{equation}
No isotropy assumption $C=I/N$ is made.

\paragraph{Theorem 1 (exact random-orientation moments).}
Conditionally on any fixed $C$,
\begin{align}
 \mathbb E_P[\rho\mid C]&=1,\\
 \operatorname{Var}_P(\rho\mid C)
 &=\frac{2(N-r)\left(N\Tr(C^2)-1\right)}
 {r(N-1)(N+2)}.
 \label{eq:rho-exact-var}
\end{align}

\paragraph{Proof.}
Rotational invariance of the Grassmannian gives $\mathbb E P=(r/N)I$, hence the first identity.  The real Grassmannian second moment for a rank-$r$ projector gives
\begin{equation}
 \operatorname{Var}_P[\Tr(PC)\mid C]
 =\frac{2r(N-r)\left(N\Tr(C^2)-1\right)}
 {N^2(N-1)(N+2)}.
\end{equation}
Multiplying by $(N/r)^2$ yields Eq.~\eqref{eq:rho-exact-var}. \hfill$\square$

\paragraph{Theorem 2 (anisotropy-independent concentration bound).}
For every admissible $N,r,C$,
\begin{equation}
 \operatorname{Var}_P(\rho\mid C)
 \le \frac{2\Tr(C^2)}{r}
 =\frac{2}{r d_{\rm eff}(C)}.
 \label{eq:rho-deff-bound}
\end{equation}
Consequently, for every $\varepsilon>0$,
\begin{equation}
 \Pr_P\!\left(|\rho-1|\ge\varepsilon\mid C\right)
 \le \min\!\left\{1,
 \frac{2}{\varepsilon^2 r d_{\rm eff}(C)}\right\}.
 \label{eq:rho-chebyshev}
\end{equation}

\paragraph{Proof.}
Write $q=\Tr(C^2)$.  Since $1/N\le q\le1$, Eq.~\eqref{eq:rho-exact-var} and
\begin{equation}
 \frac{(N-r)(Nq-1)}{(N-1)(N+2)}\le q
\end{equation}
imply Eq.~\eqref{eq:rho-deff-bound}.  Equation~\eqref{eq:rho-chebyshev} follows from Chebyshev's inequality. \hfill$\square$

\paragraph{Corollary 1 (asymptotic rank typicality).}
Consider any sequence $(N_n,r_n,C_n)$ with $N_n\to\infty$, $1\le r_n<N_n$, $C_n\succeq0$, and $\Tr C_n=1$.  If
\begin{equation}
 r_n d_{{\rm eff},n}\longrightarrow\infty,
\end{equation}
then for independent Haar/Grassmann relative orientation,
\begin{equation}
 \rho_n\xrightarrow[n\to\infty]{P}1.
\end{equation}
Thus rank typicality does not require $d_{{\rm eff},n}/N_n\to1$ and is compatible with increasingly anisotropic covariance sequences.

\paragraph{Corollary 2 (fixed-weight readout law).}
For the centered computational-basis score space of an $n$-qubit system,
\begin{equation}
 N_n=2^n-1.
\end{equation}
For diagonal Pauli-$Z$ strings of weight at most fixed $k$,
\begin{equation}
 r_k(n)=\sum_{j=1}^{k}{n\choose j}
 =\frac{n^k}{k!}\left(1+O(n^{-1})\right).
\end{equation}
Because $d_{{\rm eff},n}\ge1$, $r_k(n)d_{{\rm eff},n}\to\infty$ for every fixed $k\ge1$.  Hence under random relative orientation,
\begin{equation}
 \rho_k\xrightarrow{P}1,
\end{equation}
and therefore
\begin{equation}
 \Tr(P_kC_n)
 =(1+o_P(1))\frac{r_k(n)}{2^n-1}
 =(1+o_P(1))\frac{n^k}{k!\,2^n}.
 \label{eq:fixed-weight-exp-scaling}
\end{equation}
The accessible fraction can therefore be exponentially small in $n$ even when the global tangent covariance is strongly anisotropic.

\paragraph{Corollary 3 (almost-sure form).}
Place the sequence of random projectors on a common probability space.  If
\begin{equation}
 \sum_{n=1}^{\infty}\frac{1}{r_n d_{{\rm eff},n}}<\infty,
\end{equation}
then $\rho_n\to1$ almost surely by Eq.~\eqref{eq:rho-chebyshev} and the first Borel--Cantelli lemma.  For fixed $k\ge2$, this condition holds automatically because $d_{{\rm eff},n}\ge1$ and $r_k(n)=\Theta(n^k)$.  For $k=1$, additional growth of $d_{{\rm eff},n}$ is sufficient but not required for convergence in probability.

\subsection{Bridge to a fixed physical readout}
\label{sec:physical-orientation-bridge}

The preceding results concern a random relative orientation and do not by themselves assert that a particular physical low-weight readout is Haar-random.  They nevertheless define a scale against which a fixed physical projector can be compared.  Let
\begin{equation}
 \sigma_{\rho,n}^2
 =\frac{2(N_n-r_n)(N_n\Tr(C_n^2)-1)}
 {r_n(N_n-1)(N_n+2)}
\end{equation}
and define, when $\sigma_{\rho,n}>0$,
\begin{equation}
 z_n^{\rm(pop)}
 =\frac{\rho_n^{\rm(phys)}-1}{\sigma_{\rho,n}}.
 \label{eq:z-pop}
\end{equation}
This is an identity-based orientation diagnostic, not an additional randomization assumption.  Equation~\eqref{eq:z-pop} implies
\begin{equation}
 |\rho_n^{\rm(phys)}-1|
 =|z_n^{\rm(pop)}|\,\sigma_{\rho,n}
 \le |z_n^{\rm(pop)}|\sqrt{\frac{2}{r_nd_{{\rm eff},n}}}.
 \label{eq:physical-bridge-bound}
\end{equation}
Therefore the sufficient condition
\begin{equation}
 |z_n^{\rm(pop)}|=o\!\left(\sqrt{r_nd_{{\rm eff},n}}\right)
\end{equation}
implies $\rho_n^{\rm(phys)}\to1$.  Conversely, a symmetry can evade the rank-typical regime by generating an orientation anomaly whose growth compensates for, or exceeds, the shrinking random-orientation scale.

\subsection{Status of the numerical bridge analysis}

The rigorous-v2 primary inference was frozen before outcomes were inspected.  The population-null diagnostic in Eq.~\eqref{eq:z-pop} is introduced after inspection of those outcomes and is therefore labeled exploratory.  Its numerical implementation uses the already archived pairwise U-statistic estimator of $\Tr(C^2)$ and does not run new quantum circuits.  Rows in which the U-statistic estimate falls outside the population PSD interval $[1/N,1]$ are retained but their population-null standardized diagnostics are marked invalid rather than clipped into a favorable result.
